\section{Canonical differential polynomials and finite jets}
\label{sec:algebra}

\subsection{The normal-form convention}

Work over $\Q$ in
\[
 \Rcal=\Q[u_0,u_1,u_2,\ldots],\qquad Du_r=u_{r+1},
 \qquad \wt(u_r)=r+1.
\]
The factor degree of a monomial $\prod u_r^{a_r}$ is $\sum a_r$. A monomial is \emph{canonical} or \emph{functional} if it is a power of $u_0$, including $1$, or if its largest positive derivative index occurs at least twice. Write $\Fcal$ for the linear span of these monomials. This is the functional-monomial complement used in the classical theory of integro-differential algebras \cite[Proposition 4.10]{Guo2014} and in the repository \cite{RepoSector}.

\begin{proposition}[Canonical decomposition]\label{prop:normal}
Every $P\in\Rcal$ has a unique decomposition
\begin{equation}\label{eq:normal}
 P=\Nc(P)+D\Hc(P),\qquad \Nc(P)\in\Fcal,
 \qquad [1]\Hc(P)=0.
\end{equation}
The maps are linear, preserve factor degree, and respectively preserve and lower weight by one. Reduction never increases the largest derivative index.
\end{proposition}

\begin{proof}
It suffices first to reduce a noncanonical monomial $M$ whose unique largest index is $r\ge1$. Replace its factor $u_r$ by $u_{r-1}$, obtaining $T$, and let $a$ be the exponent of $u_{r-1}$ in $T$. In $DT$, differentiation of any of those $a$ factors gives the same monomial $M$. Every other term has largest index at most $r-1$. Consequently
\[
 M=\frac1aDT-\frac1a\bigl(DT-aM\bigr).
\]
Repeating this operation on the remaining terms terminates, because their largest derivative index strictly decreases at each nontrivial branch. Record the primitives $T/a$ to construct $\Hc$.

For uniqueness, suppose $DH\in\Fcal$ and $H$ is nonconstant. If $r$ is the largest index occurring in $H$, then $DH$ has the nonzero part
$u_{r+1}\partial H/\partial u_r$. Every monomial in this part has unique largest index $r+1$, and no other term of $DH$ can cancel it. This contradicts $DH\in\Fcal$. Hence $H$ is constant and is zero under the chosen normalization. Weight and factor-degree assertions follow directly from the reduction.
\end{proof}

This is a vector-space decomposition, not a multiplicative quotient. For example, $\Nc(u_1)=0$ but $\Nc(u_1^2)=u_1^2$. In particular one cannot multiply normal-form classes as if $D\Rcal$ were an ideal.

\subsection{Bell, Appell, and Gamma identities}

The exponential generating series for \eqref{eq:hierarchy-intro} is
\begin{equation}\label{eq:Bell}
 \Ac(t)=\sum_{m\ge0}\Ac_m\frac{t^m}{m!}
 =\exp\left(\sum_{r\ge0}\frac{(-1)^r u_rt^{r+1}}{(r+1)!}\right).
\end{equation}
Indeed $(u_0-D)\Ac(t)=\partial_t\Ac(t)$, with constant coefficient one. Thus the recurrence and the exponential define the same hierarchy.

Translation $u_0\mapsto u_0+c$ commutes with $D$ and preserves the canonical monomial space, so it commutes with $\Nc$. It follows that the normalized series has an Appell factor:
\begin{equation}\label{eq:appell}
 \sigma_d\Nc\Ac(t)=e^{u_0t}F_d(x_1,\ldots,x_d),
 \qquad x_r=u_rt^{r+1},
\end{equation}
where $\sigma_d$ sets all $u_r$, $r>d$, equal to zero. The coefficients of $F_d$ are initially formal. The entire integral formula below will prove their analytic convergence at every finite $d$.

For a smooth function $U$, put $u_r=U^{(r)}$. If $Y'=-UY$, induction gives
\begin{equation}\label{eq:YBell}
 \Ac_m(U)=(-1)^m\frac{Y^{(m)}}{Y}.
\end{equation}
Taking $U=\psi$ and $Y=1/\Gamma$ proves the Gamma specialization. This argument also prepares the Stieltjes transport without assuming that its function is itself a logarithmic derivative of Gamma.

\subsection{The variational operator}

Define
\begin{equation}\label{eq:Euler}
 \Ec(P)=\sum_{r\ge0}(-D)^r\partial_{u_r}P.
\end{equation}
This is a finite sum for each polynomial. The commutator
\[
 \partial_{u_r}D=D\partial_{u_r}+\partial_{u_{r-1}},
 \qquad \partial_{u_{-1}}=0,
\]
gives $\Ec(DH)=0$ by cancellation.

The following finite-jet statement is inherited from the prior report; we repeat its proof because the order of projection and normalization is essential.

\begin{lemma}[Finite-jet variational principle]\label{lem:finitejet}
Let $\delta_d=\sum_{r=0}^{d-1}u_{r+1}\partial_{u_r}$. Then
\begin{equation}\label{eq:finitejet}
 \sigma_d\Ec(P)=\sum_{r=0}^d(-\delta_d)^r
                 \partial_{u_r}\bigl(\sigma_d\Nc(P)\bigr).
\end{equation}
\end{lemma}
\begin{proof}
Replace $P$ by $\Nc(P)$, since their difference is a total derivative. Any canonical monomial removed by $\sigma_d$ contains its largest index greater than $d$ at least twice. A partial derivative removes at most one of these factors. Subsequent applications of $D$ leave at least one factor of index greater than $d$, so the whole contribution vanishes after projection. On retained polynomials, the ideal generated by $u_{d+1},u_{d+2},\ldots$ is $D$-stable, and projection of $D$ is $\delta_d$. This proves the formula term by term.
\end{proof}

Differentiating \eqref{eq:Bell} with respect to $u_r$ and applying \eqref{eq:Euler} gives the formal identity
\begin{equation}\label{eq:EA}
 \Ec\Ac(t)=\sum_{r\ge0}\frac{t^{r+1}}{(r+1)!}D^r\Ac(t)
          =\int_0^t e^{sD}\Ac(t)\dd s.
\end{equation}
Here integration means coefficientwise integration of a formal series. After projecting to $u_0,\ldots,u_d$, the exponent in $e^{sD}\Ac(t)$ equals
\[
 \sum_{r=0}^d\frac{u_r}{(r+1)!}
             \bigl(s^{r+1}-(s-t)^{r+1}\bigr).
\]
Substitute $s=t(1+v)/2$ to obtain
\begin{equation}\label{eq:EAcentered}
 \sigma_d\Ec\Ac(t)=\frac{te^{u_0t}}2
     \int_{-1}^1\exp\left(\sum_{r=1}^d C_r(v)x_r\right)\dd v.
\end{equation}

Set
\begin{equation}\label{eq:Ld}
 L_d=x_1+\sum_{r=1}^{d-1}x_{r+1}\partial_{x_r}.
\end{equation}
Conjugation by $e^{u_0t}$ turns $\delta_d$ into $t^{-1}L_d$. Combining \eqref{eq:finitejet} and \eqref{eq:EAcentered} gives
\begin{equation}\label{eq:general-PDE}
 \boxed{\displaystyle
 F_d+\sum_{r=1}^d(-1)^rL_d^r\partial_{x_r}F_d
 =\frac12\int_{-1}^1\exp\left(\sum_{r=1}^dC_r(v)x_r\right)\dd v.}
\end{equation}
The problem is to solve this equation under the canonical support restrictions. The next two sections solve it uniformly rather than recursively in $d$ or weight.
